\documentclass[10pt,a4paper]{amsart}
\usepackage[margin=19mm]{geometry}
\usepackage{amsmath,amssymb,mathtools}
\usepackage[scr=rsfs,cal=euler]{mathalfa}
\usepackage{xfrac}
\usepackage[unicode,hidelinks,bookmarks=false]{hyperref}
\newcommand{\ie}{i.e.\ }
\newcommand{\cf}{cf.\ }
\newcommand{\resp}{respectively\ }
\newcommand{\Subsection}{\subsection}
\providecommand{\nobreakdash}{\nobreak\hbox{-}}
\setlength{\emergencystretch}{2em}
\multlinegap0pt
\usepackage{amssymb}
\usepackage{float}
\usepackage{tikz}
\usepackage{xcolor}
\newtheorem{lemma}{Lemma}
\newtheorem{theorem}{Theorem}
\usepackage{cleveref}
\crefname{lemma}{Lemma}{Lemmas}
\crefname{theorem}{Theorem}{Theorems}
\newcommand{\R}{\mathbb{R}}
\title{Data-driven discovery and control of multistable nonlinear systems and hysteresis via structured Neural ODEs}
\author{Ike Griss Salas, Ethan King}
\date{Source: arXiv:2603.27024v1 (CC BY 4.0)}
\begin{document}
\maketitle

\noindent\emph{Source and licence.} Data-driven discovery and control of multistable nonlinear systems and hysteresis via structured Neural ODEs. Version arXiv:2603.27024v1 (CC BY 4.0), distributed by arXiv under the Creative Commons Attribution 4.0 International licence, http://creativecommons.org/licenses/by/4.0/, as recorded in the arXiv licence metadata for this exact version and shown on the abstract page https://arxiv.org/abs/2603.27024v1. Full licence notice: \texttt{licenses/arxiv-2603.27024-CC-BY-4.0.txt}.\par\smallskip

\noindent\textit{Editorial scope.} Counted unit: the article's theorem thm:local-contraction (source0.tex lines 300--310). The article's complete original proof of it is delivered with it (source0.tex lines 312--326, one \texttt{proof} environment of 15 lines). No mathematics is written, repaired or re-derived: the statement and the proof are the article's own lines, in the article's own notation, and no retained line is substituted or re-flowed.  Retained uncounted as the source-local context the proof needs: lines 275--286.  Nothing was omitted from any retained span, so the package carries no omission record.  No retained line contains a citation, so the wrapper carries no bibliography and no bibliography entry.  No retained line contains a figure environment, an image-inclusion command or drawing code, so no image is delivered and the package declares no asset dependency.  Editorial formatting only, in addition: an AMS \texttt{amsart} preamble that restates the article's own declarations and macros used by the retained lines, namely the environments \texttt{lemma}, \texttt{theorem} on separate counters, together with the standard packages those lines need.  The retained environments print in this document as Lemma 1, Theorem 1 in order of appearance, whereas the article prints the same items with the numbering of its own full document.  The complete native source of the article as distributed in the arXiv e-print for this exact version is bundled unchanged as \texttt{sources/197319/source0.tex}.
\begin{lemma}[Linear stability condition]\label{lem:linear-stability}
Let \(x^* \in \R\) be a fixed point of the scalar system
\[
    \frac{\mathrm{d}x}{\mathrm{d}t} = f(x)\bigl(x - g(x,u)\bigr),
\]
for a given constant input \(u \in \mathcal{U}\) (i.e. \(x^* = g(x^*,u)\)). Suppose \(f\in C^1(\R)\)  with \(f(x^*)<0\), and 
the partial map \(g(\cdot, u) \in C^1(\R)\).  
Then \(x^*\) is a locally exponentially stable equilibrium if and only if
\[
    \frac{\partial}{\partial x} g(x^*,u) < 1.
\]
\end{lemma}

\begin{theorem}[Local contraction of \(g\)]\label{thm:local-contraction}
Under the hypotheses of \cref{lem:linear-stability}, assume in addition that
\[
    \left|\frac{\partial}{\partial x} g(x^*,u)\right| < 1.
\]
Then there exist constants \(L\in(0,1)\) and \(r>0\) such that \(g(\cdot,u)\) is a contraction on the closed ball \(B_r(x^*)\) with Lipschitz constant \(L\), i.e.
\[
    |g(x,u) - g(y,u)| \le L |x-y|, \quad \text{for}\quad x,y\in B_r(x^*),
\]
and moreover \(g(B_r(x^*))\subset B_r(x^*)\).
\end{theorem}

\begin{proof}
Because \(g\) is \(C^1\) in \(x\) and \(|\frac{\partial}{\partial x} g(x^*,u)|<1\), continuity of \(\frac{\partial}{\partial x} g(\cdot,u)\) gives \(r>0\) and \(L\in(0,1)\) such that
\[
    \sup_{x\in B_r(x^*)} \left|\frac{\partial}{\partial x} g(x,u)\right| \le L.
\]
For any \(x,y\in B_r(x^*)\) the mean-value (integral) form yields
\[
    |g(x,u)-g(y,u)| = \left|\int_y^x \frac{\partial}{\partial x} g(s,u)\,\mathrm{d}s\right| \le L|x-y|,
\]
so \(g\) is Lipschitz on \(B_r(x^*)\) with constant \(L<1\). Taking \(y=x^*\) we obtain
\[
    |g(x,u)-x^*| \le L|x-x^*| < Lr < r,
\]
hence \(g(B_r(x^*))\subset B_r(x^*)\).
\end{proof}

\end{document}
